\section{Results and Failure Analysis}\label{sec:results}

\subsection{Recorded full-scale validation result}\label{sec:full_results}

Table~\ref{tab:full_results} reports the final scheduled monitoring evaluation
recorded for the selected full-scale slot/type run, at epoch 28,159 after
approximately 11.07 billion simulated frames. The separately downloaded
checkpoint is from epoch 28,170. These values are post-correction,
morphology-matched evaluations, not measurements of that later checkpoint
or of the independent test split.

\begin{table}[t]
\centering
\small
\caption{Recorded full-scale monitoring snapshot. Validation contains
1,048 held-out clips; each evaluation uses one selected trained body per clip.
The active resident pool is difficulty-biased. Values are the stored
distributed logging summaries from run \texttt{3mvgad6f}, epoch 28,159;
no confidence intervals or all-shape coverage are implied.}\label{tab:full_results}
\begin{tabular}{lrr}
\toprule
Metric & Hard resident pool & Validation holdout \\
\midrule
Success rate (\%) $\uparrow$ & 89.84 & 98.47 \\
Mean body-position error (m) $\downarrow$ & 0.0960 & 0.0698 \\
Mean global body-rotation error (rad) $\downarrow$ & 0.2261 & 0.1760 \\
Mean maximum-joint position error (m) $\downarrow$ & 0.1848 & 0.1412 \\
Windowed normalized jerk $\downarrow$ & 1,442.8 & 1,151.7 \\
High-jerk windows (\%) $\downarrow$ & 31.05 & 27.25 \\
\bottomrule
\end{tabular}
\end{table}

The validation score is higher than the resident-pool score, while the
reported errors are lower. This comparison is consistent with a curriculum
that concentrates residency on difficult clips: the resident pool is not an
independent uniform sample of the training distribution. Its 89.84\% success
rate should not be interpreted as the success rate of the entire training
corpus, nor should the difference be described as validation outperforming
an IID training sample.

The 98.47\% result supports strong tracking of the selected held-out
clip/body pairs under the specified threshold. It does not establish
all-body success for those clips, unseen-body generalization, or an untouched
test-set score. The thresholded success rate is complemented by position,
rotation, and smoothness measures because a successful rollout can still
have visible local errors or jitter.

% Evidence: note/README.note.md, section 76, recorded 2026-09-07.
% Do not replace the monitoring epoch with last.ckpt's epoch.
The run record is available at
\href{https://wandb.ai/yugoamaryl/hhi-protomotions/runs/3mvgad6f}
{the full-scale W\&B run}. The current table uses the recorded snapshot,
not a newly executed evaluation.
\pending{Add frozen-panel validation and independent test results from the
chosen checkpoint, with per-motion records and explicit coverage.}

\subsection{Corrected architecture comparison}\label{sec:architecture_results}

An earlier evaluator could assign a reference to an environment with the
wrong morphology. Its correction changed the meaning of the evaluation:
a rise at the code transition cannot be attributed to instantaneous policy
learning. The accompanying sampling correction also changed which training
clips entered the resident pool. These changes must be distinguished from
an architectural improvement.

The five corrected runs in Table~\ref{tab:ablations} reassess the 150-clip
comparisons using the matching-asset protocol. Their final results are not
yet included. Consequently, the selected slot/type architecture is the
current system choice, not a demonstrated winner over basic attention,
a temporal MLP, or actor-only adaLN-Zero under the corrected controls.
Pre-correction curves cannot establish a reliable architecture ranking.

\pending{For A--D, report matched-budget success, position/rotation errors,
smoothness, parameter counts, and GPU-hours on identical evaluation pairs.
Add repeated seeds before making claims about consistency.}

\subsection{Reference refinement: quality and policy effects}\label{sec:refinement_results}

Reference refinement has two distinct evaluation targets. The first is
whether the processed reference has less sliding, penetration, or jitter
without excessive changes to pose, root trajectory, or cross-shape
variation. The second is whether a policy trained on those references
tracks more reliably or smoothly. A favorable reference-only metric does
not establish the second effect.

The C/E comparison isolates the reference-processing choice at fixed
architecture and training budget, subject to common final evaluation
targets. Before/after reference metrics should be aggregated per clip/body
pair, with distributions and failure cases rather than only a favorable
mean. Changes in global geometry are expected across bodies; retained
shape-specific behavior should be assessed through matched clips and
interpretable rotation or motion-range differences.

\pending{Add the corpus-level reference artifact/fidelity audit and completed
C/E controller comparison. A full-scale refinement benefit is not inferred
causally from the single refined full-scale run.}

\subsection{Cross-shape consistency and qualitative evidence}

A useful qualitative comparison places several bodies performing the same
clip side by side, with reference markers and a shared time axis. Source
descriptions can identify the action, provided mirrored and cropped captions
have been checked. Reference-only videos illustrate the generated corpus
but must not be presented as successful policy rollouts.

For quantitative cross-shape analysis, the same clip should be evaluated
on an explicit body panel. Report both aggregate accuracy and variation
across shapes, including the number of clips that succeed on every tested
body. A rotating one-shape snapshot does not provide this statistic.
Unseen beta vectors require corresponding new assets and references and
must be evaluated separately from trained bodies, even if their coefficients
lie within the same $[-3,3]$ range.

\pending{Insert matched reference/controller visualizations and cross-shape
results. Unseen-body and joint motion/body generalization remain unmeasured
in the evidence reported here.}

\subsection{Evidence-based analysis of residual failures}\label{sec:failures}

We do not assume a fixed 10\% failure tail or an intrinsic 80\% learning
ceiling. The observed failure fraction depends on the evaluated population,
shape panel, checkpoint, horizon, and threshold. Moreover, the old
morphology mismatch produced false failures. A residual failure should
therefore be treated as an observation to explain, not evidence that the
motion is physically impossible.

The analysis unit is a failed clip/body pair with a reproducible rollout.
For each pair, retain the source identity, morphology, first threshold
crossing, tracking-error trajectory, and a synchronized reference/controller
visualization. Compare it with a successful body variant of the same clip
where available. Table~\ref{tab:failure_evidence} defines candidate
explanations and the evidence needed to support them.

\begin{table}[t]
\centering
\small
\caption{Residual-failure analysis framework. Categories are hypotheses until
supported by measurements; they can overlap and include an unresolved group.}\label{tab:failure_evidence}
\begin{tabularx}{\linewidth}{@{}p{0.25\linewidth}Y@{}}
\toprule
Candidate explanation & Evidence required \\
\midrule
Reference artifact &
Sliding, penetration, abrupt rotation/root changes, or inconsistent contacts
at the failure time; compare refined and original references. \\
Missing environmental support &
Reference/video and checked caption show an object, seat, wall, partner,
or other support absent from the simulated task. \\
Morphology-specific difficulty &
The same clip succeeds on other bodies; compare geometry, clearance,
joint limits, and reference differences on the failing body. \\
Controller or optimization limitation &
Repeated failure despite an apparently valid reference; recovery under
additional targeted training or an independently tested controller would
argue against intrinsic infeasibility. \\
Metric or protocol effect &
Inspect asset identity, horizon, initialization, and threshold crossings;
check whether a visually plausible deviation alone changes the binary label. \\
Unresolved &
Retain ambiguous and mixed cases without assigning a physical cause. \\
\bottomrule
\end{tabularx}
\end{table}

A complete report should quantify these categories over all evaluated
failures, or over a documented stratified sample with appropriate
denominators, and include representative successes as controls. Repeated
rollouts and targeted interventions can strengthen an explanation; visual
inspection alone cannot prove dynamic infeasibility. Capacity limits,
reference defects, and body-dependent dynamics may coexist.
\pending{Populate failure frequencies and representative case studies from
the frozen evaluation records. No causal failure attribution is claimed yet.}
